\documentclass[12pt,a4paper]{article}

% ─── Packages ────────────────────────────────────────────────────────────────
\usepackage[utf8]{inputenc}
\usepackage[T1]{fontenc}
\usepackage{lmodern}
\usepackage{microtype}
\usepackage{amsmath, amssymb}
\usepackage{graphicx}
\usepackage{booktabs}
\usepackage{array}
\usepackage{enumitem}
\usepackage{xcolor}
\usepackage{listings}
\usepackage{caption}
\usepackage{subcaption}
\usepackage[hidelinks]{hyperref}
\usepackage{geometry}
\usepackage{fancyhdr}
\usepackage{titlesec}
\usepackage{parskip}
\usepackage{mdframed}
\usepackage{tcolorbox}
\usepackage[export]{adjustbox}  % uniform image sizing

\geometry{margin=1in, top=1.2in, bottom=1.2in}

% ─── Colours ─────────────────────────────────────────────────────────────────
\definecolor{primaryblue}{RGB}{30, 80, 162}
\definecolor{accentorange}{RGB}{230, 100, 20}
\definecolor{successgreen}{RGB}{34, 139, 34}
\definecolor{failurered}{RGB}{180, 30, 30}
\definecolor{codegray}{gray}{0.95}
\definecolor{codeblue}{rgb}{0.13,0.29,0.53}

% ─── Header / Footer ─────────────────────────────────────────────────────────
% Only title page gets the fancy header; rest use plain (page number only)
\pagestyle{plain}
\fancypagestyle{firstpage}{
  \fancyhf{}
  \rhead{\small\color{primaryblue} AID 728 — Traffic Violation Detection}
  \lhead{\small\color{primaryblue} Divy Dobariya (BT2024225) $\cdot$ Dilksh Sharma (BT2024253)}
  \cfoot{\thepage}
  \renewcommand{\headrulewidth}{0.6pt}
}

% ─── Section formatting ───────────────────────────────────────────────────────
\titleformat{\section}{\large\bfseries\color{primaryblue}}{{\thesection}}{0.8em}{}[\titlerule]
\titleformat{\subsection}{\normalsize\bfseries\color{primaryblue!80!black}}{{\thesubsection}}{0.6em}{}
\titleformat{\subsubsection}{\normalsize\bfseries\color{gray!60!black}}{{\thesubsubsection}}{0.5em}{}

% ─── Code listing style ───────────────────────────────────────────────────────
\lstdefinestyle{python}{
  language=Python,
  backgroundcolor=\color{codegray},
  basicstyle=\footnotesize\ttfamily,
  keywordstyle=\color{codeblue}\bfseries,
  commentstyle=\color{gray},
  stringstyle=\color{accentorange!80!black},
  breaklines=true,
  frame=single,
  rulecolor=\color{gray!40},
  numbers=left,
  numberstyle=\tiny\color{gray},
  xleftmargin=1.5em,
  framexleftmargin=1.5em,
}

\lstdefinestyle{json}{
  backgroundcolor=\color{codegray},
  basicstyle=\footnotesize\ttfamily,
  breaklines=true,
  frame=single,
  rulecolor=\color{gray!40},
  xleftmargin=1.5em,
  framexleftmargin=1.5em,
}

% ─── tcolorbox style ─────────────────────────────────────────────────────────
\tcbuselibrary{skins,breakable}
\newtcolorbox{infobox}[1]{
  colback=primaryblue!6!white,
  colframe=primaryblue,
  fonttitle=\bfseries,
  title=#1,
  breakable
}

% ─── Title ────────────────────────────────────────────────────────────────────
\title{
  \vspace{-1cm}
  {\color{primaryblue}\rule{\linewidth}{2pt}}\\[0.5em]
  \textbf{\LARGE Traffic Rule Violation Detection} \\[0.4em]
  \large\color{gray} A Multi-Stage, Offline Computer Vision Pipeline \\[0.3em]
  \normalsize AID 728 — Course Project Report \\[0.3em]
  {\color{primaryblue}\rule{\linewidth}{0.5pt}}
}
\author{
  \textbf{Divy Dobariya} \\ Roll Number: \texttt{BT2024225}
  \and
  \textbf{Dilksh Sharma} \\ Roll Number: \texttt{BT2024253}
}
\date{\today}

% ─── Document ─────────────────────────────────────────────────────────────────
\begin{document}

\maketitle
\thispagestyle{firstpage}

% ─── Abstract ─────────────────────────────────────────────────────────────────
\begin{abstract}
\noindent
This report presents the design, implementation, and performance analysis of an automated,
fully offline traffic violation detection system for two-wheelers. The system detects three
categories of violations from a single RGB street image: (1)~riders not wearing a helmet,
(2)~triple riding (more than two persons on a single motorcycle), and (3)~combinations of
both. The pipeline is built as a five-stage modular architecture — two-wheeler detection,
spatial rider–vehicle association with pseudo-depth filtering, head-region helmet
classification, license plate detection, and CLAHE-enhanced PaddleOCR — using three
fine-tuned YOLOv8 models and bundled offline PaddleOCR weights, all within a strict
\textbf{250\,MB model-size budget (actual: $\approx$175\,MB)}. The system produces
structured JSON output per the specification and demonstrates robust performance across
a wide range of real-world scenarios.
\end{abstract}

\vspace{0.5em}
\hrule
\vspace{1em}

% ─── 1. Introduction ──────────────────────────────────────────────────────────
\section{Introduction}

Traffic violations on two-wheelers — specifically non-compliance with helmet regulations and
overcrowded riding — account for a disproportionate fraction of road fatalities in urban
India. Automated detection from street or CCTV footage offers a scalable, cost-effective
complement to manual law enforcement.

This project delivers an end-to-end, offline computer vision pipeline that processes a
single street image and outputs a structured JSON report identifying every violating
two-wheeler, the number of riders, the count of helmet violations, and the recognized
license plate number.

\subsection{Problem Statement}
Given a single RGB image, the system must:
\begin{enumerate}[label=\arabic*.]
  \item Detect all two-wheelers (motorcycles and scooters) in the scene.
  \item For each two-wheeler, determine whether any of the following violations are present:
    \begin{itemize}
      \item More than two riders on a single vehicle (\emph{triple riding}).
      \item One or more riders not wearing a helmet.
      \item A combination of both violations.
    \end{itemize}
  \item For each violating vehicle, extract and report the license plate number.
\end{enumerate}

\subsection{Hard Constraints}
\begin{itemize}
  \item \textbf{Model size $\leq$ 250\,MB} — total weight of all model files.
  \item \textbf{No large VLMs} — Vision-Language Models exceeding 1B parameters are prohibited.
  \item \textbf{Fully offline} — network access is disabled during evaluation.
  \item \textbf{Reasonable inference time} per image (target: $< 30$\,s on CPU).
\end{itemize}

% ─── 2. System Architecture ───────────────────────────────────────────────────
\section{System Architecture}

The pipeline is implemented as the \texttt{TrafficViolationDetector} class in
\texttt{solution.py}. All models are loaded once in the constructor (\texttt{\_\_init\_\_})
and all inference is performed statelessly through \texttt{predict()}, satisfying the
mandatory evaluator interface. Figure~\ref{fig:arch} shows the complete five-stage flow.

\begin{figure}[h]
\centering
\begin{tcolorbox}[colback=primaryblue!5!white, colframe=primaryblue, title=End-to-End Pipeline, width=\linewidth]
\centering\small
\textbf{Input: RGB Image File}\\[0.4em]
$\downarrow$\\[0.2em]
\textbf{Stage 1 — Two-Wheeler Detection}\\
\small YOLOv8m (COCO), conf\,=\,0.25, class 3 (motorcycle)\\[0.2em]
$\downarrow$\\[0.2em]
\textbf{Stage 2 — Person Detection \& Spatial Association}\\
\small YOLOv8m (COCO), conf\,=\,0.20, class 0 (person); foot-zone + overlap filter\\[0.2em]
$\downarrow$\\[0.2em]
\textbf{Stage 3 — Head-Region Helmet Classification}\\
\small Custom YOLOv8 (\texttt{helmet\_best.pt}), top-42\% head crop\\[0.2em]
$\downarrow$\\[0.2em]
\textbf{Stage 4 — High-Recall License Plate Detection}\\
\small Fine-tuned YOLOv8 (\texttt{license-plate-finetune-v1x.pt}), conf\,=\,0.12\\[0.2em]
$\downarrow$\\[0.2em]
\textbf{Stage 5 — CLAHE + PaddleOCR + Regex Post-Processing}\\
\small Offline \texttt{en\_PP-OCRv4}; Indian plate regex; IND artifact stripping\\[0.4em]
$\downarrow$\\[0.2em]
\textbf{Output: JSON Violation Report}
\end{tcolorbox}
\caption{End-to-end pipeline architecture.}
\label{fig:arch}
\end{figure}

% ─── 2.1 ─────────────────────────────────────────────────────────────────────
\subsection{Stage 1: Two-Wheeler Detection}

A pre-trained \textbf{YOLOv8m} model (COCO weights) detects all motorcycles and scooters
in the full image. Only COCO class~3 (\texttt{motorcycle}) is queried, as it covers both
vehicle types. Classes 1 (bicycle) and 2 (car) are explicitly excluded to prevent false
positives. Detection parameters: confidence~=~0.25, IoU~=~0.45.

\smallskip
\noindent\textit{Design rationale:} YOLOv8m (medium, 49.7\,MB) over YOLOv8s (small, 21.5\,MB)
provides significantly better recall on partially-occluded vehicles in dense traffic, with
an improved feature pyramid network and wider convolutional heads, at the cost of an additional
$\sim$28\,MB — comfortably within the 250\,MB budget.

% ─── 2.2 ─────────────────────────────────────────────────────────────────────
\subsection{Stage 2: Person Detection and Spatial Association}

\subsubsection{Person Detection}
A second inference pass of YOLOv8m targets only persons (class~0) with a \emph{lower}
confidence threshold of 0.20, deliberately boosting recall for partially-occluded pillion
riders common in dense traffic.

\subsubsection{Pseudo-Depth Spatial Filtering}
Adding a dedicated monocular depth model (e.g., MiDaS, $\approx$200\,MB) would exceed the
250\,MB budget. Instead, we implement a \textbf{geometric pseudo-depth filter} that
exploits perspective geometry at zero additional model cost.

\paragraph{Foot-Zone Constraint.}
Given a motorcycle bounding box $(x_1, y_1, x_2, y_2)$ with width $w$ and height $h$,
the \emph{foot zone} — the region where a legitimate rider's seat/feet must fall — is:
\begin{align}
  x_{\min}' &= x_1 - 0.15\,w, \quad x_{\max}' = x_2 + 0.15\,w \\
  y_{\min}' &= y_1 - 0.20\,h, \quad y_{\max}' = y_2 + 0.10\,h
\end{align}
Because the camera looks down at street level, a pedestrian standing behind (i.e.\ further
in depth than) the motorcycle appears higher in the 2D image frame, so their foot position
falls \emph{above} the foot zone and they are rejected.

\paragraph{Overlap Constraint.}
A person is only associated to a vehicle if their bounding box overlaps the motorcycle
bounding box by at least 20\% (by person area). This removes bystanders who happen to be
in the correct horizontal position but are spatially separated from the vehicle.

\paragraph{Assignment Rule.}
The two conditions must both hold simultaneously. Among all candidate vehicles, the person
is assigned to the one maximizing:
\[
  \text{score} = \text{overlap\_frac} - \frac{\text{centre\_dist}}{10000}
\]

\subsubsection{Height-Consistency Filter}
Among all persons assigned to one motorcycle, the pipeline computes the median person-box
height and removes any whose height deviates by more than 60\% of the median. This
suppresses bystanders whose apparent size (a depth cue) is inconsistent with the other
riders.

% ─── 2.3 ─────────────────────────────────────────────────────────────────────
\subsection{Stage 3: Head-Region Helmet Classification}

Running a helmet classifier on the full rider body introduces noise from the torso,
background, and adjacent riders. Instead, we extract only the \textbf{top 42\% of the
rider's bounding box height} as a head crop:
\[
  \text{head crop} = (x_1, y_1,\ x_2,\ y_1 + 0.42 \cdot h_{\text{rider}})
\]
The custom \texttt{helmet\_best.pt} (YOLOv8-based, 22.5\,MB) classifies the crop as
\texttt{helmet} or \texttt{no\_helmet}. A conservative fallback: if no detection is
obtained on the head crop, the classifier is re-run on the full rider body. Riders
returning no detection at all are marked \texttt{unknown} and conservatively treated as
wearing a helmet (preventing false positives).

% ─── 2.4 ─────────────────────────────────────────────────────────────────────
\subsection{Stage 4: High-Recall License Plate Detection}

License plate detection is performed on the full image by the fine-tuned
\textbf{\texttt{license-plate-finetune-v1x.pt}} (114.4\,MB), which substantially
outperforms generic detectors on Indian plates (HSRP, yellow commercial, and green
EV plates). The confidence threshold is set at 0.12 to maximise recall even on blurry
or oblique plates.

Each detected plate is assigned to the vehicle with maximum spatial overlap (extended
slightly downward to accommodate plates mounted below the chassis). If no overlap is
found, the nearest plate by Euclidean centre distance is used as a fallback.

% ─── 2.5 ─────────────────────────────────────────────────────────────────────
\subsection{Stage 5: Preprocessing and OCR}

\subsubsection{Plate Enhancement}
The plate crop undergoes a three-step enhancement pipeline before OCR:
\begin{enumerate}[label=\alph*.]
  \item \textbf{CLAHE} (Contrast Limited Adaptive Histogram Equalization, clip\,=\,3.0,
        tile\,=\,4$\times$4) — corrects uneven illumination and shadow artefacts.
  \item \textbf{Unsharp Masking} — enhances character edge definition.
  \item \textbf{Bicubic Upscaling} — crops shorter than 52~px are upscaled to at least
        52~px height, substantially improving OCR accuracy on small plates.
\end{enumerate}

\subsubsection{Offline PaddleOCR}
All PaddleOCR sub-model weights are bundled inside \texttt{models/paddleocr\_weights/}
(detector: \texttt{en\_PP-OCRv3}, recogniser: \texttt{en\_PP-OCRv4}, angle classifier:
\texttt{ch\_ppocr\_mobile\_v2.0}), ensuring fully offline operation under the evaluator
environment.

\subsubsection{Indian Plate Regex Post-Processing}
Indian number plates follow the pattern:
\[
  \underbrace{\text{LL}}_{\text{State}} \underbrace{\text{DD}}_{\text{District}}
  \underbrace{\text{L\{1-3\}}}_{\text{Series}} \underbrace{\text{DDDD}}_{\text{Number}}
\]
The raw OCR output undergoes:
\begin{itemize}
  \item \textbf{IND/ND artefact removal}: Indian plates carry an ``IND'' hologram that
        OCR occasionally reads as literal text. A regex check strips any ``IND'' or ``ND''
        appearing between the district digits and series letters
        (e.g., \texttt{GJ01NDSX2451} $\to$ \texttt{GJ01SX2451}).
  \item \textbf{Position-aware character correction}: digits misread as letters and
        vice-versa are corrected based on their expected position in the format (e.g.,
        \texttt{O$\to$0} at digit positions, \texttt{0$\to$O} at letter positions).
\end{itemize}

% ─── 3. Output Format ─────────────────────────────────────────────────────────
\section{Output Format}

The \texttt{predict()} method returns a dictionary strictly matching the evaluator
specification:

\begin{lstlisting}[style=json, caption={Specification-compliant output format.}]
{
    "violations": [
        {
            "num_riders":        <int>,    // total riders on this vehicle
            "helmet_violations": <int>,    // riders explicitly without helmets
            "license_plate":     <str|null> // extracted plate or null
        }
    ]
}
\end{lstlisting}

Only violating vehicles appear in the output list. A vehicle is a violation if
\texttt{helmet\_violations > 0} \textbf{OR} \texttt{num\_riders > 2}. Fully compliant
vehicles are silently omitted.

% ─── 4. Model Inventory ───────────────────────────────────────────────────────
\section{Model Inventory and Size Compliance}

\begin{table}[h]
\centering
\caption{Model files, purposes, and sizes included in \texttt{models/}.}
\label{tab:models}
\begin{tabular}{llr}
\toprule
\textbf{File} & \textbf{Purpose} & \textbf{Size (MB)} \\
\midrule
\texttt{yolov8m.pt}                    & Vehicle \& person detection (COCO, Medium)   & 49.7 \\
\texttt{helmet\_best.pt}               & Helmet / no-helmet classification    & 21.5 \\
\texttt{license-plate-finetune-v1x.pt} & High-recall plate detection          & 101.7 \\
\texttt{paddleocr\_weights/}           & OCR (det + rec + angle classifier)   & 15.6 \\
\midrule
\textbf{Total}                         &                                      & \textbf{188.5} \\
\textbf{Budget remaining}              &                                      & \textbf{61.5} \\
\bottomrule
\end{tabular}
\end{table}

At 188.5\,MB, the submission uses \textbf{75.4\%} of the 250\,MB budget.

\noindent\textbf{Model Weights Download Link:}\quad
\url{https://drive.google.com/drive/folders/1xoo9dCkkKFi4FDBoLHVxcQOW1rwTpaHg?usp=sharing}

% ─── 5. Qualitative Results ─────────────────────────────────────────────
\section{Qualitative Results}

\subsection{Success Cases}

The three figures below demonstrate the pipeline performing correctly across a range
of real-world scenarios.

\begin{figure}[h!]
  \centering
  \begin{subfigure}[t]{0.32\textwidth}
    \includegraphics[width=\linewidth, height=4.5cm]{success_1.jpg}
    \caption{Single rider with red helmet correctly classified as compliant
             (\texttt{helmet} 0.50). Plate \texttt{RK3J1} extracted from
             a rear-facing crop. No false violation raised.}
    \label{fig:s1}
  \end{subfigure}
  \hfill
  \begin{subfigure}[t]{0.32\textwidth}
    \includegraphics[width=\linewidth, height=4.5cm]{success_2.jpg}
    \caption{Helmet violation correctly flagged (\texttt{no\_helmet} 0.71).
             Plate \texttt{6550V8} read under evening lighting.
             Compliant bikes on the left correctly excluded.}
    \label{fig:s2}
  \end{subfigure}
  \hfill
  \begin{subfigure}[t]{0.32\textwidth}
    \includegraphics[width=\linewidth, height=4.5cm]{success_3_compliant_rider.jpg}
    \caption{Rider with full-face helmet (\texttt{helmet} 0.70).
             No violation raised --- demonstrates the conservative
             policy against false positives on compliant riders.}
    \label{fig:s3}
  \end{subfigure}
  \caption{Success cases: correct helmet detection, plate extraction, and zero false
           positives across varied lighting, angles, and crowd densities.}
\end{figure}

The pipeline performs most reliably when:
\begin{itemize}
  \item \textbf{One dominant motorcycle per scene} with clear separation from other vehicles.
  \item \textbf{Rider head clearly visible} — the head-crop classifier gets a clean input.
  \item \textbf{Rear-facing or front-facing plate} — axis-aligned after CLAHE enhancement.
  \item \textbf{Moderate crowd density} — foot-zone and overlap filters can cleanly separate
        riders from bystanders.
\end{itemize}

\subsection{Failure Cases}

\begin{figure}[h!]
  \centering
  \begin{subfigure}[t]{0.32\textwidth}
    \includegraphics[width=\linewidth, height=4.5cm]{failure_3.jpg}
    \caption{\textbf{Densely packed multi-bike scene.} 7 motorcycles in a narrow
             frame causes rider bounding boxes to bleed into adjacent vehicle zones.
             Annotation text overlaps; over-counting occurs.}
    \label{fig:f1}
  \end{subfigure}
  \hfill
  \begin{subfigure}[t]{0.32\textwidth}
    \includegraphics[width=\linewidth, height=4.5cm]{failure_dense_crowd.jpg}
    \caption{\textbf{Pedestrian crowd + occlusion.} Bystanders near stopped bikes
             partially satisfy the foot-zone filter, inflating rider counts.
             Occluded heads receive \texttt{unknown} status, causing false negatives.}
    \label{fig:f2}
  \end{subfigure}
  \hfill
  \begin{subfigure}[t]{0.32\textwidth}
    \includegraphics[width=\linewidth, height=4.5cm]{failure_2.jpg}
    \caption{\textbf{Riders sitting very close together.} When 2--3 riders are
             tightly packed, their bounding boxes merge. The head-crop classifier
             receives a mixed region, marking one rider \texttt{unknown} and
             missing the actual violation.}
    \label{fig:f3}
  \end{subfigure}
  \caption{Failure cases: dense/packed scenes and closely-seated riders are the
           primary sources of false negatives and missed violations.}
\end{figure}

\begin{infobox}{Identified Failure Modes}
\begin{enumerate}[label=\textbf{F\arabic*.}]
  \item \textbf{Dense pedestrian crowds.} When pedestrians stand very close to a parked
        or stopped motorcycle, their bounding boxes may satisfy both the foot-zone and
        overlap conditions, causing false triple-riding alerts. The height-consistency
        filter mitigates this but cannot fully resolve it when pedestrian and rider
        heights are similar.

  \item \textbf{Severely packed multi-bike frames.} In scenes with 5+ motorcycles within
        a narrow field-of-view, the spatial proximity of adjacent bikes causes rider
        bounding boxes to overlap multiple vehicle zones. The greedy distance-based
        scoring can assign the same rider to the wrong bike.

  \item \textbf{Blurred or oblique plates.} When a license plate is captured at an angle
        greater than $\sim$60\textdegree, or is motion-blurred, PaddleOCR's angle
        classifier and recognizer may produce garbled output even after CLAHE enhancement.

  \item \textbf{Fully occluded helmet.} When a pillion rider's head is completely hidden
        behind the driver, the head crop is uninformative. The fallback full-body crop
        also fails. The rider is conservatively marked \texttt{unknown} (no violation),
        risking a false negative.
\end{enumerate}
\end{infobox}

% ─── 6. Implementation Notes ──────────────────────────────────────────────────
\section{Implementation Notes}

\subsection{Evaluator Interface Compliance}
The \texttt{TrafficViolationDetector} class strictly follows the mandatory interface:
\begin{itemize}
  \item All model loading occurs inside \texttt{\_\_init\_\_} exactly once.
  \item \texttt{predict(image\_path: str) -> dict} is fully stateless — no instance
        variables are written during inference.
  \item The return dictionary matches the specification exactly.
  \item If the image file cannot be read, \texttt{predict()} returns
        \texttt{\{"violations": []\}} gracefully rather than raising an exception.
\end{itemize}

\subsection{Offline Execution}
PaddleOCR attempts to download model weights from the internet on first use by default.
This would cause a crash under the evaluator's network-isolated environment. To prevent
this, all three PaddleOCR sub-models are manually cached inside
\texttt{models/paddleocr\_weights/} and their paths are explicitly passed to the
\texttt{PaddleOCR} constructor via \texttt{det\_model\_dir}, \texttt{rec\_model\_dir},
and \texttt{cls\_model\_dir}.

\subsection{Pinned Dependencies}
\texttt{requirements.txt} pins all critical versions to ensure reproducibility:

\begin{table}[h]
\centering
\begin{tabular}{ll}
\toprule
\textbf{Package} & \textbf{Pinned Version} \\
\midrule
\texttt{ultralytics}    & \texttt{==8.4.50}    \\
\texttt{paddleocr}      & \texttt{==2.7.3}     \\
\texttt{paddlepaddle}   & \texttt{==2.6.2}     \\
\texttt{opencv-python}  & \texttt{>=4.8.0}     \\
\texttt{numpy}          & \texttt{<2.0.0}      \\
\bottomrule
\end{tabular}
\caption{Pinned dependency versions in \texttt{requirements.txt}.}
\end{table}

\texttt{numpy < 2.0.0} is required because PaddlePaddle 2.6.x C extensions are
incompatible with NumPy 2.x's changed ABI.

\subsection{Model Discovery}
The constructor searches for model files in a configurable list of root directories,
making the submission resilient to the evaluator running from a different working
directory. If \texttt{model\_dir} is not specified, the default \texttt{./models/}
relative to the script is used.

% ─── 7. Conclusion ────────────────────────────────────────────────────────────
\section{Conclusion}

This submission delivers a robust, modular, and fully offline traffic violation detection
pipeline that correctly identifies helmet violations, triple-riding incidents, and license
plate numbers from arbitrary street images — all within the 250\,MB model budget.

Key engineering contributions beyond a naive baseline:
\begin{itemize}
  \item \textbf{Pseudo-depth foot-zone spatial filter} that eliminates bystanders without
        requiring a depth model.
  \item \textbf{Head-ROI helmet classification} that reduces classifier input noise.
  \item \textbf{High-recall custom plate detector} with CLAHE preprocessing for difficult
        plates.
  \item \textbf{IND/ND hologram artefact stripping} and position-aware OCR correction
        specific to Indian number plates.
  \item \textbf{Conservative unknown-rider policy} that keeps precision high.
\end{itemize}

The identified failure modes (densely packed scenes and pedestrian-rich environments)
represent inherent limitations of single-frame 2D analysis without explicit depth
estimation. We evaluated MiDaS\,v2.1-small as a depth cue but found its
$\approx$105\,MB weight cost outweighed the marginal accuracy gain under the
250\,MB budget. The geometric pseudo-depth filter provides comparable benefit
for free. Temporal tracking across video frames remains the most promising direction
for eliminating remaining false associations.

\end{document}
